\chapter{Padrões de Projeto e UML}
\label{cap:padroes-uml-teoria}

\begin{edital}
“Padrões de projeto”, “design de software”, “padrões de desenvolvimento e
reuso”, “análise e projeto orientados a objetos”. UML não está listada com
todas as letras no edital do Perfil 3 (aparece explícita no Perfil 2), mas é
vizinha direta desses itens e a FGV cobra com frequência em provas de TI.
\end{edital}

\section{Padrões de projeto (GoF)}

\begin{conceito}[O que é --- e o que não é --- um padrão de projeto]
Um padrão de projeto \textbf{não é código pronto, nem biblioteca, nem
framework}. É a descrição de uma solução que já foi reinventada tantas
vezes, por gente diferente, para o mesmo problema estrutural, que virou
vocabulário compartilhado. Como todo verbete de dicionário técnico, tem
quatro partes: um \textbf{nome} (para poder falar sobre ele sem reexplicar),
o \textbf{problema} em que se aplica, a \textbf{solução} em termos de
classes e objetos que colaboram, e as \textbf{consequências} --- porque toda
escolha de design tem um preço (mais classes, mais indireção, mais
complexidade para quem lê o código pela primeira vez).

Essa definição já entrega o principal erro de quem estuda padrão decorando
a estrutura: aplicar um padrão \textbf{sem ter o problema} correspondente é
\textit{over-engineering}. Uma \textit{Abstract Factory} para criar um
único tipo de objeto, sem variação nenhuma, acrescenta duas classes e zero
flexibilidade real --- o padrão existe para resolver um problema, não para
demonstrar conhecimento de padrões.

O catálogo GoF (\textit{Gang of Four}, os quatro autores do livro que o
consolidou em 1994) reúne \textbf{23 padrões}, divididos em três famílias:
\textbf{criacionais} (5), \textbf{estruturais} (7), \textbf{comportamentais}
(11).
\end{conceito}

\begin{conceito}[O fio que liga quase todos os 23]
Por trás da variedade, a maioria dos padrões do catálogo faz a
\textbf{mesma} manobra estrutural: \textbf{isola o que varia atrás de uma
interface estável}, e prefere \textbf{composição a herança} para poder
trocar essa parte variável em tempo de execução, sem recompilar nem alterar
quem usa a interface. É o fio que liga Strategy, State, Bridge, Decorator,
Abstract Factory e a maior parte dos outros --- e é por isso que “programe
para uma interface, não para uma implementação” e “prefira composição a
herança” são os dois lemas mais citados do livro original.

O que muda de família para família é \textbf{o que} exatamente está sendo
isolado:

\begin{center}
\begin{tabular}{@{}p{3.2cm}p{3.8cm}p{4.8cm}@{}}
\toprule
\textbf{Família} & \textbf{Isola\ldots} & \textbf{Pergunta que ela responde} \\
\midrule
Criacionais (5) & \textit{qual} objeto se instancia, e quando & “como criar
isto sem amarrar o código à classe concreta?” \\
Estruturais (7) & \textit{como} as peças se encaixam & “como
compor/embrulhar objetos para obter uma interface ou um comportamento
novo?” \\
Comportamentais (11) & \textit{quem} conversa com quem, e qual algoritmo
roda & “como distribuir a responsabilidade da interação em tempo de
execução?” \\
\bottomrule
\end{tabular}
\end{center}
\end{conceito}

\subsection{Criacionais}

\begin{conceito}
\begin{center}
\begin{tabular}{@{}p{2.8cm}p{9cm}@{}}
\toprule
\textbf{Padrão} & \textbf{Intenção (gatilho no enunciado)} \\
\midrule
Singleton & garantir \textbf{uma única instância} e ponto global de acesso
(“apenas uma conexão/configuração no sistema todo”) \\
Factory Method & delegar a criação a \textbf{subclasses}; criar objeto
\textbf{sem acoplar à classe concreta} (“decidir em runtime qual objeto
criar”, via herança) \\
Abstract Factory & criar \textbf{famílias} inteiras de objetos relacionados
sem especificar classes concretas (“kit de componentes de UI para Windows e
para Mac”) \\
Builder & construir objeto complexo \textbf{passo a passo}, separando o
processo de construção da representação final (“montar um objeto com
muitos parâmetros opcionais”) \\
Prototype & criar novos objetos \textbf{clonando} um protótipo existente
(“copiar um objeto já configurado em vez de criar do zero”) \\
\bottomrule
\end{tabular}
\end{center}
\end{conceito}

\begin{exemplo}[Factory Method x Abstract Factory]
Um sistema de notificação precisa enviar mensagens por e-mail ou SMS,
decidido em tempo de execução por configuração. Uma classe abstrata
\texttt{NotificadorFactory} tem o método \texttt{criar()}, e cada subclasse
(\texttt{EmailFactory}, \texttt{SmsFactory}) decide qual notificador
concreto instanciar --- isso é \textbf{Factory Method}: uma família de
objetos, delegada a subclasses por herança.

Agora imagine que o sistema de UI precisa gerar um \textbf{conjunto
inteiro} de componentes visualmente consistentes: botão, campo de texto e
janela, todos no estilo “Windows” ou todos no estilo “Mac”, nunca
misturados. Uma \texttt{UIFactory} abstrata declara \texttt{criarBotao()},
\texttt{criarCampo()}, \texttt{criarJanela()}; \texttt{WindowsFactory} e
\texttt{MacFactory} implementam a família inteira de uma vez --- isso é
\textbf{Abstract Factory}: não um objeto, uma \textbf{família} de objetos
relacionados que precisam ser consistentes entre si.
\end{exemplo}

\begin{cuidado}
Factory Method resolve \textbf{um} produto por vez, via herança (uma
subclasse por variante). Abstract Factory resolve uma \textbf{família}
inteira de produtos relacionados, via composição (a fábrica concreta é
injetada, não subclassificada por quem a usa). Singleton é sempre
\textbf{uma} instância --- nunca “mais de uma quando necessário”, isso
descaracteriza o padrão.
\end{cuidado}

\subsection{Estruturais}

\begin{conceito}
\begin{center}
\begin{tabular}{@{}p{2.8cm}p{9cm}@{}}
\toprule
\textbf{Padrão} & \textbf{Intenção (gatilho)} \\
\midrule
Adapter & fazer interfaces \textbf{incompatíveis} trabalharem juntas
(“adaptar uma API legada à interface nova que o sistema espera”) \\
Facade & interface \textbf{simplificada} para um subsistema complexo (“uma
fachada única escondendo vários serviços internos”) \\
Proxy & um \textbf{substituto} que controla o acesso ao objeto real
(carregamento tardio, cache, controle de permissão) \\
Decorator & adicionar responsabilidades a um objeto \textbf{dinamicamente},
sem usar herança (“empilhar comportamentos em tempo de execução”) \\
Composite & tratar objeto individual e composição de objetos \textbf{de
forma uniforme} (estrutura em árvore, relação todo-parte recursiva) \\
Bridge & separar \textbf{abstração} de \textbf{implementação} para as duas
variarem de forma independente \\
Flyweight & compartilhar objetos semelhantes para \textbf{economizar
memória} (muitas instâncias quase idênticas) \\
\bottomrule
\end{tabular}
\end{center}
\end{conceito}

\begin{exemplo}[Adapter x Facade x Decorator x Proxy --- os quatro que se confundem]
Um sistema novo espera receber pagamentos por uma interface
\texttt{ProcessadorPagamento.pagar(valor)}, mas a biblioteca do banco só
oferece \texttt{BancoLegado.executarTransacao(valorEmCentavos, moeda)}.
Uma classe \texttt{AdapterBanco} implementa \texttt{ProcessadorPagamento} e,
por dentro, converte a chamada para o formato do banco legado --- isso é
\textbf{Adapter}: faz duas interfaces que já existem e não combinam
conversarem.

O mesmo sistema, para fechar um pedido, precisa chamar \texttt{Estoque},
\texttt{Pagamento}, \texttt{NotaFiscal} e \texttt{Envio}, em sequência
determinada. Uma classe \texttt{FachadaPedido.finalizar(pedido)} esconde
essas quatro chamadas atrás de um método só --- isso é \textbf{Facade}:
simplifica o acesso a um subsistema que já funciona, sem mudar nenhuma das
peças internas.

Agora um sistema de café precisa somar o preço base a extras escolhidos
(leite, chantilly, calda) sem criar uma subclasse para cada combinação
possível. \texttt{CafeComLeite} envolve um \texttt{Cafe} e soma seu próprio
preço ao preço de quem envolve; \texttt{CafeComChantilly} faz o mesmo em
volta do resultado anterior --- isso é \textbf{Decorator}: cada camada
\textbf{acrescenta comportamento/dado}, empilhável em qualquer combinação.

Por fim, uma imagem de 50MB só deve ser carregada do disco quando
realmente for exibida na tela. Uma classe \texttt{ProxyImagem} implementa a
mesma interface de \texttt{Imagem}, mas só instancia a \texttt{ImagemReal}
(que faz o carregamento pesado) na primeira chamada a \texttt{exibir()}
--- isso é \textbf{Proxy}: \textbf{controla o acesso} ao objeto real, sem
acrescentar comportamento novo à funcionalidade em si.
\end{exemplo}

\begin{cuidado}[A diferença que sobra depois da estrutura parecida]
Adapter, Facade e Decorator todos “envolvem” outro objeto, e por isso a
estrutura de código fica parecida. O que os separa é a \textbf{intenção}:
Adapter \textbf{traduz} uma interface existente para outra que o cliente
espera (o objetivo é compatibilidade); Facade \textbf{simplifica} o acesso
a várias peças que já são compatíveis entre si (o objetivo é conveniência);
Decorator \textbf{acrescenta} funcionalidade nova, mantendo a mesma
interface (o objetivo é extensão). E Proxy \textbf{controla acesso} sem
mudar a função --- é a diferença que separa Proxy de Decorator: o primeiro
decide \textit{se e quando} delega para o objeto real; o segundo sempre
delega, e soma algo à resposta.
\end{cuidado}

\subsection{Comportamentais}

\begin{conceito}
\begin{center}
\begin{tabular}{@{}p{2.8cm}p{9cm}@{}}
\toprule
\textbf{Padrão} & \textbf{Intenção (gatilho)} \\
\midrule
Strategy & família de algoritmos \textbf{intercambiáveis}, escolhidos por
\textbf{quem usa}, em runtime (“trocar a regra de cálculo sem um
\texttt{if} gigante”) \\
Observer & notificar automaticamente \textbf{dependentes} quando o estado
de um objeto muda (\textit{publish-subscribe}, sistema de eventos) \\
Command & encapsular uma \textbf{requisição como objeto} (permite
desfazer/refazer, enfileirar, logar operações) \\
State & mudar o comportamento conforme o \textbf{estado interno}, e é o
\textbf{próprio objeto/estado} quem decide a próxima transição \\
Template Method & define o \textbf{esqueleto} fixo de um algoritmo,
deixando passos específicos para subclasses sobrescreverem \\
Iterator & percorre uma coleção \textbf{sem expor sua estrutura interna} \\
Chain of Responsibility & passa a requisição por uma \textbf{cadeia} de
manipuladores até que algum a trate \\
Mediator & centraliza a comunicação entre vários objetos, reduzindo
acoplamento \textbf{N-para-N} para N-para-1 \\
Memento & captura e restaura o estado interno de um objeto sem violar seu
encapsulamento (\textit{snapshot}, base do “desfazer”) \\
Visitor & adiciona novas operações a uma estrutura de objetos \textbf{sem
alterar as classes} dessa estrutura \\
Interpreter & define uma \textbf{gramática} para uma linguagem simples e um
interpretador que avalia sentenças escritas nela \\
\bottomrule
\end{tabular}
\end{center}
\end{conceito}

\begin{exemplo}[Do \texttt{if} gigante ao Strategy]
O cálculo de frete começa assim, dentro do serviço de pedidos:
\texttt{if (tipo == "sedex") \{...\} else if (tipo == "pac") \{...\} else
if (tipo == "retirada") \{...\}}. Refatorando com Strategy: cria-se uma
interface \texttt{CalculoFrete} com o método \texttt{calcular(pedido)}, e
uma classe por regra (\texttt{FreteSedex}, \texttt{FretePac},
\texttt{FreteRetirada}). O serviço passa a \textbf{receber} a estratégia já
escolhida e só chamar \texttt{calcular} --- ele deixa de conhecer as regras
concretas, dependendo só da interface.

O ganho tem nome: acrescentar uma transportadora nova virou \textbf{criar
uma classe}, não \textbf{editar} um método já existente --- é o
\textit{Open/Closed} do SOLID (Seção~\ref{sec:solid-uml} adiante) em ação.
E o serviço passou a depender de uma abstração (\texttt{CalculoFrete}), não
de implementações concretas --- é o \textit{Dependency Inversion} do mesmo
princípio. Repare que o padrão não eliminou a decisão de \textbf{qual}
regra usar; ele só a \textbf{moveu} para fora, para quem monta o objeto do
serviço.
\end{exemplo}

\begin{cuidado}[Strategy x State --- mesma estrutura, intenção oposta]
Os dois são, por dentro, a mesma estrutura de código: um objeto delegando
comportamento a outro através de uma interface comum. A diferença está em
\textbf{quem decide a troca}. No \textbf{Strategy}, quem escolhe a
estratégia é \textbf{o cliente}, de fora, e as estratégias não sabem umas
das outras (trocar o cálculo de frete não afeta o cálculo de imposto). No
\textbf{State}, quem troca é \textbf{o próprio objeto}, conforme seu estado
interno, e os estados costumam conhecer para qual estado ir a seguir
(“pedido pago” sabe que o próximo é “pedido enviado”). Regra de leitura: se
o enunciado descreve um \textbf{ciclo de vida com transições} (pedido
pendente $\to$ pago $\to$ enviado $\to$ entregue), é State; se descreve
\textbf{alternativas equivalentes} para fazer a mesma tarefa (calcular
frete de três jeitos diferentes), é Strategy.
\end{cuidado}

\begin{cuidado}[Observer x Mediator]
\textbf{Observer} é o padrão de notificação: um sujeito muda de estado e
\textbf{avisa} uma lista de interessados que se inscreveram para saber
disso (a comunicação é sempre “sujeito $\to$ observadores”, um-para-muitos).
\textbf{Mediator} \textbf{centraliza a comunicação} entre vários objetos
que, de outra forma, teriam que se conhecer uns aos outros (N-para-N vira
N-para-1 através do mediador) --- não é sobre notificar mudança de estado,
é sobre reduzir quantas classes conhecem quantas outras classes
diretamente.
\end{cuidado}

\subsection{GRASP}

\begin{conceito}[Princípios de responsabilidade, não padrões prontos]
GRASP (\textit{General Responsibility Assignment Software Patterns}) não é
uma lista de estruturas de código como o GoF --- é um conjunto de
\textbf{princípios} para decidir \textbf{qual classe deve receber qual
responsabilidade}: Information Expert, Creator, Controller, Low Coupling,
High Cohesion, Polymorphism, Pure Fabrication, Indirection, Protected
Variations.

Dois deles resolvem sozinhos a maior parte das perguntas de “a qual classe
cabe esta responsabilidade”:

\begin{itemize}
\item \textbf{Information Expert}: a responsabilidade vai para a classe
que \textbf{já tem os dados} necessários para cumpri-la. Quem calcula o
total de um pedido é a classe \texttt{Pedido}, que já conhece seus itens
--- não uma \texttt{CalculadoraDeTotais} externa, que precisaria pedir
tudo emprestado (o que criaria acoplamento desnecessário).
\item \textbf{Creator}: quem cria um objeto é quem o \textbf{agrega,
contém ou usa de perto}. \texttt{Pedido} cria seus \texttt{ItemDePedido},
porque é quem os contém e usa diretamente.
\end{itemize}

Os dois existem pelo mesmo motivo: manter juntos o dado e o comportamento
que o usa é o que produz \textbf{alta coesão} (a classe faz uma coisa
coerente) e \textbf{baixo acoplamento} (menos classes precisam se conhecer
para o trabalho acontecer) --- que são, eles próprios, dois outros
princípios da mesma lista.
\end{conceito}

\subsection{SOLID}
\label{sec:solid-uml}

\begin{conceito}[Os cinco princípios --- tratamento completo com armadilhas de Java no Capítulo~\ref{cap:java-teoria}]
\begin{itemize}
\item \textbf{S --- \textit{Single Responsibility}}: uma classe deve ter
\textbf{uma responsabilidade}, ou seja, um único motivo para mudar. Uma
classe que salva no banco \textbf{e} formata relatório \textbf{e} valida
regra de negócio tem três motivos para mudar --- e muda toda vez que
qualquer um dos três muda, mesmo sem relação com os outros dois.
\item \textbf{O --- \textit{Open/Closed}}: aberta à \textbf{extensão},
fechada à \textbf{modificação}. Acrescentar comportamento novo deve ser
possível \textbf{sem editar} código já testado e em produção --- é
exatamente o que o exemplo de Strategy acima demonstrou.
\item \textbf{L --- \textit{Liskov}}: um subtipo deve poder substituir seu
tipo base em qualquer lugar do programa \textbf{sem quebrar o
comportamento esperado}. O exemplo clássico de violação: \texttt{Quadrado}
herda de \texttt{Retangulo} e sobrescreve \texttt{setLargura} para também
mudar a altura (mantendo os lados iguais) --- um código que usa
\texttt{Retangulo} esperando poder mudar largura e altura independentemente
quebra quando recebe um \texttt{Quadrado}.
\item \textbf{I --- \textit{Interface Segregation}}: várias interfaces
\textbf{pequenas e específicas} são melhores que uma interface “gorda” que
obriga quem implementa a fornecer métodos que não usa.
\item \textbf{D --- \textit{Dependency Inversion}}: módulos de alto nível
não devem depender de módulos de baixo nível --- ambos devem depender de
\textbf{abstrações}. Na prática: o serviço depende da interface
\texttt{CalculoFrete}, nunca da classe concreta \texttt{FreteSedex}.
\end{itemize}
\end{conceito}

\begin{cuidado}[ISP x SRP]
Os dois falam em “pequeno e específico”, e por isso se confundem, mas
tratam de coisas diferentes: \textbf{SRP} é sobre quantas
\textbf{responsabilidades} uma \textbf{classe} tem; \textbf{ISP} é sobre
quantos \textbf{métodos desnecessários} uma \textbf{interface} obriga a
implementar. Uma classe pode violar SRP mesmo implementando uma interface
enxuta, e uma interface pode violar ISP mesmo sendo implementada por uma
classe de responsabilidade única.
\end{cuidado}

\subsection{Padrões arquiteturais (fora do catálogo GoF)}

\begin{conceito}
\begin{itemize}
\item \textbf{MVC} (\textit{Model-View-Controller}): separa dados
(\textit{Model}), apresentação (\textit{View}) e controle do fluxo
(\textit{Controller}) em três responsabilidades distintas.
\item \textbf{MVVM / MVP}: variações do MVC com \textit{binding} de dados
automático (MVVM) ou um \textit{presenter} que intermedeia view e model
(MVP) --- comuns em front-end e mobile.
\item \textbf{DAO} (\textit{Data Access Object}) \textbf{/ Repository}:
isolam o código de acesso a dados (SQL, ORM) do resto da aplicação.
\item \textbf{DTO} (\textit{Data Transfer Object}): objeto simples, sem
lógica de negócio, usado só para transportar dados entre camadas.
\end{itemize}
\end{conceito}

\subsection{Reuso e anti-padrões}

\begin{conceito}
\textbf{Reuso por herança} (“é-um”: \texttt{Cachorro extends Animal}) x
\textbf{reuso por composição} (“tem-um”: \texttt{Carro} tem um
\texttt{Motor}). A recomendação consolidada é \textbf{preferir composição a
herança} (\textit{favor composition over inheritance}) sempre que a
relação não for genuinamente um tipo-subtipo estável, porque herança
acopla fortemente a subclasse aos detalhes internos da superclasse
(qualquer mudança na base pode quebrar a subclasse sem aviso).

\textbf{Anti-padrões} (más práticas recorrentes, batizadas do mesmo jeito
que os padrões, para poder nomear o erro): \textbf{God Object} (uma classe
que faz tudo, violando SRP em escala máxima), \textbf{Spaghetti Code}
(fluxo de controle emaranhado, sem estrutura reconhecível),
\textbf{Golden Hammer} (usar sempre a mesma ferramenta/solução para
problemas diferentes, mesmo quando não é a mais adequada),
\textbf{Copy-Paste Programming} (duplicar código em vez de extrair uma
abstração comum).
\end{conceito}

\section{UML}
\label{sec:uml-teoria}

\begin{conceito}[Notação, não método]
A UML é uma \textbf{linguagem de notação}, não um processo: ela padroniza
\textbf{o desenho}, não diz \textbf{quando} desenhar cada diagrama --- isso
é papel do processo de desenvolvimento (RUP, Scrum; Capítulo
\ref{cap:eng-software-teoria}). O problema que ela resolve é concreto:
antes da padronização, cada empresa tinha sua própria simbologia, e o
mesmo losango podia significar coisas diferentes em dois documentos da
mesma equipe. Padronizar a notação é o que permite discutir projeto sem
primeiro precisar combinar o que cada símbolo quer dizer. A especificação
vigente da OMG é a \textbf{UML 2.5.1}.
\end{conceito}

\subsection{Os 14 diagramas em duas famílias}

\begin{conceito}
\begin{center}
\begin{tabular}{@{}p{2.6cm}p{2.6cm}p{6.4cm}@{}}
\toprule
\textbf{Família} & \textbf{Foca em} & \textbf{Diagramas principais} \\
\midrule
Estruturais & estrutura \textbf{estática} --- como o sistema é organizado
num instante & \textbf{Classe}, Objeto, Componente, Implantação
(\textit{deployment}), Pacote, Estrutura composta, Perfil \\
Comportamentais & comportamento \textbf{dinâmico} --- como o sistema se
move no tempo & \textbf{Caso de uso}, \textbf{Sequência},
\textbf{Atividade}, \textbf{Estado} (máquina de estados), Comunicação,
Interação geral, Tempo \\
\bottomrule
\end{tabular}
\end{center}
A pergunta-teste: o diagrama descreveria o sistema \textbf{parado, num
instante} (estrutural) ou descreve \textbf{o que acontece ao longo do
tempo} (comportamental)? Classe é sempre estrutural; sequência, atividade,
estado e caso de uso são sempre comportamentais.
\end{conceito}

\subsection{Diagrama de classes}

\begin{conceito}[Relacionamentos e o que cada notação significa]
\begin{center}
\begin{tabular}{@{}p{3.6cm}p{4.6cm}p{2.6cm}@{}}
\toprule
\textbf{Relacionamento} & \textbf{Significado} & \textbf{Notação} \\
\midrule
Associação & ligação estrutural entre classes & linha \\
Agregação & todo-parte \textbf{fraca} (a parte sobrevive sem o todo) &
losango \textbf{vazio} \\
Composição & todo-parte \textbf{forte} (a parte morre junto com o todo) &
losango \textbf{cheio} \\
Herança / Generalização & é-um (subtipo do tipo base) & seta de triângulo
vazio \\
Dependência & usa temporariamente (parâmetro, variável local) & seta
tracejada \\
Realização & implementa uma interface & linha tracejada + triângulo vazio \\
\bottomrule
\end{tabular}
\end{center}

\textbf{Multiplicidade} (quantas instâncias do outro lado participam):
\texttt{1} (exatamente uma), \texttt{0..1} (zero ou uma), \texttt{1..*}
(uma ou mais), \texttt{*} (zero ou mais). \textbf{Visibilidade}:
\texttt{+} público, \texttt{-} privado, \texttt{\#} protegido,
\texttt{\textasciitilde} pacote.

A classe se desenha em \textbf{três compartimentos}: nome, atributos,
operações. O atributo, na forma completa, é onde mora o detalhe mais
cobrado em leitura de diagrama:

\begin{center}
\texttt{- status : String [1] = "aberto"}
\end{center}

Nessa ordem: visibilidade, nome, tipo, multiplicidade, e \textbf{valor
padrão} --- o valor que a instância assume quando ninguém informa outro na
criação.
\end{conceito}

\begin{exemplo}[Agregação x composição]
Uma \texttt{Universidade} tem \texttt{Departamentos} --- se a universidade
fechar, os departamentos, como entidades daquela universidade, deixam de
existir junto (não fazem sentido soltos): \textbf{composição}, losango
cheio do lado de \texttt{Universidade}. Um \texttt{Departamento} tem
\texttt{Professores} --- se o departamento fechar, os professores continuam
existindo (podem ser realocados para outro departamento):
\textbf{agregação}, losango vazio do lado de \texttt{Departamento}.
\end{exemplo}

\begin{cuidado}[Extremidade de associação como atributo, e associação qualificada]
Cada ponta de uma associação tem papel, multiplicidade e navegabilidade
--- e uma ponta navegável \textbf{pode ser representada como atributo} da
classe do outro lado. “\texttt{Pedido} associa-se a 1 \texttt{Cliente}”
desenhado como linha e “\texttt{Pedido} tem o atributo \texttt{cliente:
Cliente}” escrito no compartimento de atributos dizem a \textbf{mesma}
coisa em UML --- são duas notações de um único modelo, não duas
modelagens diferentes. Uma \textbf{associação qualificada} (um qualificador
desenhado como um pequeno retângulo na ponta da linha, geralmente uma
chave) \textbf{particiona} o conjunto do outro lado e costuma reduzir a
multiplicidade de \texttt{*} para \texttt{0..1}: dado um \texttt{Banco} e
um \texttt{numeroConta}, chega-se a no máximo uma \texttt{Conta} --- é um
índice de busca, não uma relação todo-parte, o que a torna um distrator
plausível de agregação em leitura de diagrama.
\end{cuidado}

\subsection{Casos de uso}

\begin{conceito}
Modela \textbf{requisitos funcionais} do ponto de vista de quem usa o
sistema: \textbf{ator} (bonequinho, alguém/algo externo que interage),
\textbf{caso de uso} (elipse, uma funcionalidade completa do ponto de vista
do ator), \textbf{fronteira do sistema} (retângulo que delimita o que está
dentro).

\begin{itemize}
\item \textbf{«include»}: o comportamento incluído é \textbf{sempre}
executado como parte do caso de uso base (obrigatório) --- ex.: “Fazer
Pedido” \textit{include} “Autenticar Usuário”.
\item \textbf{«extend»}: o comportamento é \textbf{opcional/condicional},
só ocorre em certas circunstâncias --- ex.: “Fazer Pedido” pode ser
\textit{extended} por “Aplicar Cupom de Desconto”, que só roda se o
cliente tiver um cupom.
\end{itemize}
\end{conceito}

\begin{cuidado}
Include é o comportamento que \textbf{sempre} acontece; extend é o que
\textbf{às vezes} acontece. É comum a leitura inverter os dois --- o teste
mental é: “esse passo é parte obrigatória do fluxo principal, ou é um
desvio condicional que só ocorre em certas situações?”
\end{cuidado}

\subsection{Os quatro diagramas de interação}

\begin{conceito}[Todos mostram objetos trocando mensagens; a diferença é a ênfase]
\begin{center}
\begin{tabular}{@{}p{3.2cm}p{4.2cm}p{4.2cm}@{}}
\toprule
\textbf{Diagrama} & \textbf{Enfatiza} & \textbf{Como se lê a ordem} \\
\midrule
\textbf{Sequência} & a \textbf{ordem cronológica} das mensagens & pela
posição vertical: de cima para baixo \\
\textbf{Comunicação} & a \textbf{estrutura de ligações} entre os objetos
(quem está conectado a quem) & pela \textbf{numeração} das mensagens (1,
1.1, 1.2) --- sem eixo de tempo explícito \\
\textbf{Tempo} (\textit{timing}) & a \textbf{mudança de estado} de cada
participante ao longo de uma \textbf{escala de tempo explícita} & pelo eixo
horizontal de tempo; usado quando o requisito envolve duração/tempo real \\
\textbf{Visão geral de interação} & como \textbf{vários cenários} se
encadeiam entre si & é um diagrama de atividades cujos nós são, eles
mesmos, interações \\
\bottomrule
\end{tabular}
\end{center}

\textbf{Sequência} e \textbf{comunicação} são \textbf{semanticamente
equivalentes} --- dá para converter um no outro sem perder informação, já
que ambos mostram as mesmas mensagens entre os mesmos objetos. A diferença
é puramente de ênfase visual: sequência é melhor para ver \textbf{quando}
cada coisa acontece; comunicação é melhor para ver \textbf{quem fala com
quem}. Diagrama de \textbf{atividade} é diferente dos quatro: ali não há
troca de mensagens entre participantes nomeados, há um \textbf{fluxo de
trabalho} genérico --- ações, decisões, paralelismo.
\end{conceito}

\begin{regrapratica}
Se o enunciado pede a \textbf{ordem exata em que as mensagens acontecem},
é \textbf{sequência}. Se pede \textbf{quem está ligado a quem}
estruturalmente, sem se importar com quando, é \textbf{comunicação}. Se
envolve \textbf{duração/tempo real}, é \textbf{timing}. Se é sobre um
\textbf{fluxo de trabalho} sem participantes trocando mensagem nomeada, é
\textbf{atividade}, não um diagrama de interação.
\end{regrapratica}

\subsection{Atividade e máquina de estados}

\begin{conceito}
\textbf{Atividade}: fluxo de trabalho/algoritmo --- nós de ação, decisão
(losango), bifurcação/junção (barra, representa paralelismo), início/fim.
Parecido com um fluxograma; usado para modelar processos de negócio ou a
lógica de um algoritmo.

\textbf{Máquina de estados}: modela os \textbf{estados possíveis} de um
objeto e as \textbf{transições} entre eles, disparadas por eventos --- o
objeto se comporta diferente conforme o estado em que está. É o diagrama
que casa diretamente com o padrão de projeto \textbf{State} (Seção
anterior): cada estado do diagrama pode virar uma classe concreta que
implementa a interface do estado.
\end{conceito}

\begin{regrapratica}[Roteiro para ligar diagrama a conceito]
Caso de uso $\leftrightarrow$ requisitos funcionais (Capítulo
\ref{cap:eng-software-teoria}). Classes $\leftrightarrow$ orientação a
objetos/SOLID. Máquina de estados $\leftrightarrow$ padrão de projeto
State. Componente e Implantação (\textit{deployment})
$\leftrightarrow$ decisões de arquitetura (Capítulo
\ref{cap:arquitetura-teoria}).
\end{regrapratica}
